\section{Physical Grounding of Discrete Information Acquisition}\label{sec:physical-grounding}

\subsection{Continuous vs Discrete: A Question of Tractability, Not Validity}\label{sec:continuous-discrete}

Both continuous and discrete models are valid. The choice between them is not about which is ``true'' but about which is tractable for the computational system using it.

\begin{itemize}
\item \textbf{Continuous models} use limits to collapse computational complexity. You don't check all states; you abstract them away. This is a useful thinking tool for mathematical analysis.

\item \textbf{Discrete models} are the actual structure that physical computational systems must work with.
\end{itemize}

\noindent The two perspectives are duals: the continuous model is the limit of the discrete model as resolution increases; the discrete model is the physical instantiation that any bounded system must use.

\noindent\textbf{Question:} Does the continuous abstraction beat physical limits?

No. Limited abstract counting collapses computational complexity in the model, but it does not beat the physical limits of computation speed. No physical decision circuit can compute faster than the speed of light. The continuous abstraction is a cognitive tool for tractable thinking; it does not help a physical computational system that is bounded by $c$.

\subsection{First-Principles Derivation: The Counting Gap}\label{sec:first-principles}

We formalize this with a theorem of pure mathematics. No physical law is assumed.

\begin{theorem}[Counting Gap Theorem]\label{thm:counting-gap}
Let $\varepsilon, C \in \mathbb{N}$ with $\varepsilon > 0$ and $C > 0$. If each information-acquisition event consumes $\varepsilon$ discrete cost units, then for every $N \in \mathbb{N}$:
\[
  \varepsilon \cdot N \leq C \;\Rightarrow\; N \leq C.
\]
Equivalently, $C$ itself is a finite upper bound on the number of possible events. No bounded system with positive integer event cost can perform infinitely many events.

\begin{proof}
This is a discrete theorem. In $\mathbb{N}$, every positive integer is $\geq 1$. Therefore $N = 1 \cdot N \leq \varepsilon \cdot N \leq C$.
\end{proof}
\end{theorem}
\leanmeta{\LH{BA10}}

\noindent\textbf{Interpretation.} A continuous process would require infinitely many events. The theorem proves that any bounded computational system with positive event cost has a finite maximum. Physical computational systems are bounded and have positive cost. Therefore physical computational systems cannot use the continuous abstraction to beat their limits.

\begin{corollary}[Forced Finiteness of Propagation Speed]\label{cor:forced-finite-speed}
A bounded region cannot acquire information at unbounded rate. Therefore some finite propagation bound $c_{\max}$ must exist.
\LH{BA10}
\end{corollary}

\subsection{Three Empirical Laws: The Numerical Content}\label{sec:three-laws}

Theorem~\ref{thm:counting-gap} proves that a finite bound exists. It does not supply the numerical value. Three empirical laws supply the values in our universe:

\begin{enumerate}
\item[\textbf{SR.}] \textbf{Special Relativity.}
  No physical circuit computes faster than $c$.
  Einstein established: $c = 2.998 \times 10^8$ m/s.
  \cite{einstein1905electrodynamics,minkowski1908space}

\item[\textbf{QM.}] \textbf{Quantum Mechanics.}
  State transitions are discrete eigenstate jumps.
  Planck and Dirac established: $E = h\nu$.
  \cite{planck1901distribution,dirac1930principles}

\item[\textbf{TD.}] \textbf{Thermodynamics (Landauer).}
  Each event costs positive energy.
  Landauer established the universal floor: $\varepsilon \geq k_B T \ln 2$ joules per irreversible bit. The Wolpert-style interface used in this artifact adds explicit nonnegative mismatch and residual dissipation terms above that floor; the mismatch term can now be justified inside the artifact from finite-distribution KL divergence when the actual and designed distributions differ, and the finite discrete residual branch can also be justified inside the artifact by an exhaustive reverse-edge split. If the reverse edge flow is positive, asymmetric forward/reverse flows yield a positive two-point KL witness; if the reverse edge flow is zero, the process performs an irreversible one-way state transition, and the existing Landauer-scaled transition-cost theorem yields a positive lower-bound witness. The broader stopping-time / absolute-irreversibility residual regime remains a separate imported premise. When any one of these contributions is positive, the per-bit lower bound is strictly larger than $k_B T \ln 2$.
  \cite{landauer1961irreversibility,bennett1982thermodynamics,berut2012experimental,wolpert2024stochastic,manzano2024absolute,yadav2026minimal}
\end{enumerate}
\leanmeta{\LH{WC1},\LH{WC2},\LH{WC3},\LH{WM1},\LH{WM3},\LH{WM4},\LH{WR1},\LH{WR2},\LH{WR3},\LH{WR6}.}

\noindent The Counting Gap Theorem proves the structure. SR, QM, TD supply the numbers.

\subsection{Why Positive Cost is Derived, Not Assumed}\label{sec:positive-cost-derived}

The premise $\varepsilon > 0$ is not arbitrary. A check is a physical state transition. Physical transitions require force. Force over distance requires energy. At $T > 0$, any irreversible transition carries at least the Landauer floor $k_B T \ln 2$. In the explicit constrained-process interface, any extra dissipation above that floor is represented by separate nonnegative mismatch and residual terms rather than folded into the floor itself. The mismatch branch is now pinned to the existing KL layer: explicit finite-distribution mismatch yields a strictly positive real KL term, which is then conservatively rounded upward into the nat-valued lower-bound units used by the thermodynamic accounting. The finite discrete residual branch is now also theorem-level by an exhaustive local split on the reverse edge of a computational-state process: either the reverse edge is positive and asymmetry yields strictly positive two-point KL divergence, or the reverse edge is zero and the process pays strictly positive Landauer-scaled cost for the resulting irreversible one-way state transition.

If $\varepsilon = 0$, a bounded system could perform infinitely many checks, resolving any decision in zero time and zero energy. This contradicts energy conservation. Therefore $\varepsilon > 0$ is derived from thermodynamics.
\leanmeta{\LH{BA10},\LH{WC2},\LH{WC3},\LH{WM1},\LH{WM3},\LH{WM4},\LH{WR2},\LH{WR3},\LH{WR6}}

\begin{remark}[Scope]
No claim is made about whether the universe is continuous or discrete at scales
below the physical resolution of any bounded system. Whether finer structure
exists is irrelevant: it is inaccessible to any bounded region of spacetime at
rate $> c/d$, and therefore has no effect on any physical decision process.
\end{remark}

\subsection{Bounded Information Acquisition Rate}\label{sec:bounded-acquisition}

\begin{proposition}[Bounded Region]\label{prop:bounded-region}
A \emph{bounded region} $\mathcal{R}$ is characterized by diameter $d > 0$ and
signal propagation speed $c > 0$ (in abstract natural units). The maximum
information-acquisition rate of $\mathcal{R}$ is $c/d$ events per unit time.

By \textbf{SR}: information emitted from outside $\mathcal{R}$ takes $\geq d/c$
time to reach $\mathcal{R}$. Therefore $\mathcal{R}$ cannot acquire information
at rate $> c/d$.
\LH{BA1}
\end{proposition}
\leanmeta{\LH{BA1}}
\claimstamp{D\ref{def:physical-core-encoding}}{AR}

\begin{theorem}[Bounded Acquisition Rate]\label{thm:bounded-acquisition}
For bounded region $\mathcal{R}$ with diameter $d$ and signal speed $c$,
operating for time $T$:
\[
\text{acquisitions}(\mathcal{R}, T) \leq \frac{c \cdot T}{d}.
\]
This count is finite. $\mathcal{R}$ cannot acquire information continuously.
\LH{BA2}
\end{theorem}
\leanmeta{\LH{BA2}}

\subsection{Acquisitions are Discrete State Transitions}\label{sec:discrete-transitions}

\begin{theorem}[Discrete Acquisition]\label{thm:discrete-acquisition}
By \textbf{QM}: the state space of any bounded region with finite energy is
countable (discrete energy eigenstates). Each information-acquisition event is
a transition $A \to B$ between eigenstates, a discrete event. Between
transitions, the state is fixed and no new information is held.

Formally: any physical decision process in $\mathcal{R}$ is a
\texttt{DiscreteSystem} (Definition in \texttt{Physics/DiscreteSpacetime.lean}).
Acquisition events are exactly \texttt{isTransitionPoint} events.
\LH{BA3}
\end{theorem}
\leanmeta{\LH{BA3}}
\claimstamp{D\ref{def:physical-core-encoding};P\ref{prop:physical-claim-transport}}{AR}

\subsection{One Transition = One Bit}\label{sec:one-bit}

\begin{theorem}[Boolean Encoding is Physically Primitive]\label{thm:boolean-primitive}
Each discrete state transition transfers exactly one boolean bit:
the system was in state $A$ (encoded $0$), it transitions to state $B$ (encoded $1$).
The encoding $\{A, B\} \cong \{0, 1\}$ is not a modeling choice; it is the
form of physical state transitions under \textbf{QM}.

Formally: each \texttt{isTransitionPoint} contributes exactly $1$ to
\texttt{bitOperations}. Boolean coordinates are the elementary quanta of
physical information exchange.
\LH{BA4}
\end{theorem}
\leanmeta{\LH{BA4}}

\begin{corollary}[Finite Accessible State Space]\label{cor:finite-state}
A bounded region $\mathcal{R}$ operating for finite time $T$ acquires at most
$n_{\max}(\mathcal{R}, T) = \lfloor cT/d \rfloor$ bits. Its accessible state
space is a finite subset of $\{0,1\}^{n_{\max}}$.
Uncountable state spaces are not accessible to any bounded region
in finite time.
\LH{BA2}\LH{BA4}
\end{corollary}
\leanmeta{\LH{BA4}, \LH{BA2}}

\subsection{Structural Rank = Minimum Physical Bit Operations}\label{sec:srank-physical}

\begin{theorem}[Resolution Requires a Sufficient Coordinate Set]\label{thm:resolution-sufficient}
Any physical process that correctly determines the optimal action for all states
of decision problem $\mathcal{D}$ must access a sufficient coordinate set $I$.

If it accesses fewer coordinates than needed for sufficiency, there exist two
states with identical readings but different optimal actions. The resolver
outputs the wrong action on one. Correctness forces sufficiency.

Each coordinate in $I$ requires one bit operation to read (one state transition:
register samples coordinate $i$ $\to$ one bit acquired by Theorem~\ref{thm:boolean-primitive}).
Therefore: bit operations $\geq |I|$.
\LH{BA5}
\end{theorem}
\leanmeta{\LH{BA5}}

\begin{theorem}[Structural Rank = Minimum Bit Operations]\label{thm:srank-physical}
For decision problem $\mathcal{D}$ with $\mathrm{srank}(\mathcal{D}) = k$:

\begin{enumerate}
\item \textbf{Information (BA6):}
  $k = \mathrm{srank}(\mathcal{D}) \leq |I|$ for any sufficient coordinate set $I$,
  by \texttt{srank\_le\_sufficient\_card} (Tractability/StructuralRank.lean).
  Therefore: minimum bits to resolve $\mathcal{D}$ correctly $= k$.

\item \textbf{Time (BA6 + SR):}
  Each bit operation requires one tick of duration $\geq d/c$.
  Minimum decision time $= k \cdot d/c$.

\item \textbf{Energy (BA7 + TD):}
  Each bit operation carries at least the Landauer floor $k_B T \ln 2$ by \textbf{TD}.
  Therefore every resolver obeys the universal lower bound $E \geq k \cdot k_B T \ln 2$; stronger empirical lower bounds only raise this floor.
\end{enumerate}

Structural rank $k$ is simultaneously an information cost, a time cost,
and an energy cost. These are not three separate quantities; they are the
same count $k$ expressed in three physical units.
\LH{BA6}\LH{BA7}
\end{theorem}
\leanmeta{\LH{BA6}, \LH{BA7}, \LH{SR1}}
\claimstamp{T\ref{thm:physical-bridge-bundle}}{AR}

\subsection{Thermodynamic Lower Bound as Physical Consequence}\label{sec:thermo-derived}

\begin{theorem}[Thermodynamic Cost Derived from SR, QM, TD]\label{thm:thermo-derived}
For any physical system resolving decision problem $\mathcal{D}$ at temperature $T > 0$:
\[
E_{\mathrm{decision}} \geq \mathrm{srank}(\mathcal{D}) \cdot k_B T \ln 2.
\]
\begin{proof}
Resolving $\mathcal{D}$ requires reading a sufficient coordinate set $I$
(Theorem~\ref{thm:resolution-sufficient}). Each coordinate read is one bit
operation (Theorem~\ref{thm:boolean-primitive}). Bit operations $\geq |I| \geq \mathrm{srank}$
(Theorem~\ref{thm:srank-physical}). By \textbf{TD}, each bit operation costs
$\geq k_B T \ln 2$. Summing: $E \geq \mathrm{srank} \cdot k_B T \ln 2$.
\end{proof}
The constant $k_B T \ln 2$ is the universal Landauer floor for irreversible
bit erasure, experimentally verified to $\pm 10\%$
\cite{berut2012experimental}. It is not, in general, a tight expression for the
actual energetic cost of constrained computation
\cite{wolpert2024stochastic,manzano2024absolute,yadav2026minimal}. In the Lean
model, that distinction is explicit in two layers. First, a coarse
floor-plus-overhead interface records that any constrained process may dissipate
more than the Landauer floor (\LH{WC1}--\LH{WC5}). Second, a refined Wolpert
decomposition splits that extra dissipation into a mismatch term and a residual
term. Theorem \LH{WP1} states that the declared total overhead is exactly the
sum of those two components, and \LH{WP2} proves that the effective per-bit
lower bound dominates the Landauer floor plus both terms. The mismatch branch
is now partially promoted from cited premise to theorem: \LH{WM1} shows the
KL mismatch term is nonnegative, \LH{WM2} identifies the zero case exactly,
\LH{WM3} gives strict positivity under an explicit finite-distribution mismatch
witness, and \LH{WM4} carries that strict positivity into the discrete lower-bound
units used by the thermodynamic model. Theorems \LH{WM5} and \LH{WM6} then
inject that theorem-level mismatch witness into the Wolpert decomposition, so
strict separation above the Landauer floor can now be derived directly for the
mismatch branch. A second residual branch is now also theorem-level within the
current finite machinery: \LH{WR1} proves nonnegativity of the two-point
forward/reverse residual divergence, \LH{WR2} proves strict positivity under an
explicit asymmetric positive edge-flow witness, and \LH{WR3} lifts that witness
into the discrete lower-bound units of the decomposition. The new theorem
\LH{WR6} packages the exhaustive finite-state edge split directly: if the
reverse edge is positive, the pairwise KL branch applies; if the reverse edge
vanishes, the process performs an irreversible one-way state transition and the
existing Landauer-scaled transition-cost theorem applies. Theorems \LH{WR7} and
\LH{WR8} then inject that unified finite residual witness into the Wolpert
decomposition, yielding direct strict separation above the Landauer floor for
that finite discrete subclass. Theorems \LH{WR9} and \LH{WR10} package the
same branch at the process level: any finite computational-state process with
at least one such asymmetric forward edge already carries a theorem-level
residual witness and therefore strictly exceeds the Landauer floor. The broader
stopping-time / absolute-
irreversibility residual regime imported from \cite{manzano2024absolute}
remains represented by the separate cited premise \LH{WP5}; and \LH{WP6} shows
that any one of the derived mismatch branch, the derived finite residual
branch, or the cited broader stopping-time branch is already sufficient to
force strict separation above the Landauer floor. For the circuit cost interface of
\cite{yadav2026minimal}, \LH{WP7} and \LH{WP8} prove that a declared structural
resource lower-bounded by mismatch further tightens the per-bit and total energy
lower bounds. Finally, \LH{WP9} proves that the existing physical grounding
bundle survives under this refined decomposition.
Exact equality is retained only as an optional idealized specialization.

\noindent\textbf{Informally:} the cited stochastic-thermodynamics papers are
not rederived in full inside this artifact. The finite-distribution mismatch
branch is now proved as far as the existing KL machinery reaches, and a finite
residual asymmetry branch is also proved from the same machinery applied to
forward/reverse edge-flow distributions, then repackaged as a single
process-level witness theorem for finite computational-state systems. The stronger stopping-time /
absolute-irreversibility residual regime remains a separate cited premise. Lean
then verifies the exact downstream consequences of composing that mixed
theorem-plus-premise interface with the decision-theoretic lower bounds already
proved in the core system.
\LH{BA9}\LH{WM1}\LH{WM2}\LH{WM3}\LH{WM4}\LH{WM5}\LH{WM6}\LH{WR1}\LH{WR2}\LH{WR3}\LH{WR4}\LH{WR5}\LH{WR6}\LH{WR7}\LH{WR8}\LH{WR9}\LH{WR10}\LH{WP5}\LH{WP6}
\end{theorem}
\leanmeta{\LH{BA9}}
\claimstamp{T\ref{thm:physical-bridge-bundle}}{AR}

\begin{corollary}[srank $= 1$ is the Energy Ground State]\label{cor:ground-state}
Among all decision problems, those with $\mathrm{srank} = 1$ have the smallest
universal lower bound certified by the structural argument:
$E_{\mathrm{decision}} \geq k_B T \ln 2$ per decision. In the exact-calibration
idealization this is the corresponding ground-state cost.

Problems with $\mathrm{srank} > 1$ are above that universal floor by at least
$(\mathrm{srank} - 1) \cdot k_B T \ln 2$ per decision cycle. This gap is
structurally mandatory at the level of lower bounds; stronger empirical
per-bit costs only enlarge it.
\LH{BA8}
\end{corollary}
\leanmeta{\LH{BA8}}
\claimstamp{C\ref{cor:ground-state}}{AR}

\subsection{Connection to the Formal Model}\label{sec:physical-to-formal}

The coordinate structure of Definition~\ref{def:decision-problem} is justified
by Theorems~\ref{thm:boolean-primitive} and~\ref{thm:srank-physical}:

\begin{itemize}
\item The finite coordinate spaces $X_1, \ldots, X_n$ correspond to the
  discrete state transitions accessible to the physical system
  (Corollary~\ref{cor:finite-state}).

\item The boolean case $X_i = \{0,1\}$ is the elementary case:
  each coordinate is one physical state transition
  (Theorem~\ref{thm:boolean-primitive}).

\item $\mathrm{srank}(\mathcal{D})$ counts the minimum physical transitions
  required to resolve $\mathcal{D}$ correctly
  (Theorem~\ref{thm:srank-physical}).

\item The thermodynamic lower bound (Theorem~\ref{thm:thermo-derived})
  is the physical content of the thermodynamic lift results
  (ThermodynamicLift.lean, DiscreteSpacetime.lean).
\end{itemize}

The formal model is substrate-neutral (Section~\ref{sec:foundations}).
The physical grounding here establishes that the model accurately describes
any physical substrate subject to \textbf{SR}, \textbf{QM}, and \textbf{TD},
which is all matter in this universe under current empirical knowledge.
\leanmeta{\LH{BA9}, \LH{DS5}}

\subsection{Universe Membership via Shared Invariants}\label{sec:universe-membership}

The phrase ``this universe'' is not rhetorical. It has precise operational content:
two physical systems are in the same universe \emph{iff} they agree on a shared
set of measurable invariants within measurement precision.

\begin{definition}[Shared Invariant Set]\label{def:invariant-set}
An \emph{invariant set} $\mathcal{I} = \{c, \hbar, k_B, \ldots\}$ is a collection
of physical constants that are:
\begin{enumerate}
\item Measurable by any sufficiently capable observer.
\item Identical (within precision) for all observers in the same universe.
\end{enumerate}
\LH{IA0}
\end{definition}
\leanmeta{\LH{IA0}}

\begin{theorem}[Universe Membership is Defined by Invariant Agreement]\label{thm:universe-membership}
Two observers $O_1, O_2$ are in the same physical universe \emph{iff} their
measurements of a shared invariant set $\mathcal{I}$ are compatible within
measurement precision.

\textbf{Contrapositive:} Without a shared measurable invariant, the phrase
``same universe'' has no operational content.
\LH{IA15}
\end{theorem}
\leanmeta{\LH{IA15}, \LH{IA16}}

This theorem is not about humans. It is grounded at the quantum level where
no ego exists, then lifted to observers as a corollary:

\begin{enumerate}
\item \textbf{Quantum level (IA2--IA4):} Two quantum systems under the same
  Hamiltonian agree on all eigenvalues. There is no ``electron perspective.''
  Agreement is forced by physics.

\item \textbf{Atomic level (IA5--IA8):} Two atoms absorbing the same photon
  agree on $\hbar\omega$. Discrete transitions force exact agreement.

\item \textbf{Molecular level (IA9--IA12):} Two molecules in thermal equilibrium
  at temperature $T$ agree on $k_B T$. The Landauer floor then applies
  universally to irreversible bit erasure.

\item \textbf{Observer level (IA13--IA18):} Observers are made of atoms and
  molecules. They inherit invariant agreement from their constituents.
  By the observer level, agreement is already forced.
\end{enumerate}
\leanmeta{\LH{IA2}, \LH{IA5}, \LH{IA9}, \LH{IA13}, \LH{IA17}, \LH{IA18}}

\begin{corollary}[Observer-Level Rejection Cascade]\label{cor:observer-rejection-cascade}
Objecting to observer-level invariant agreement requires rejecting atomic-level
invariant agreement, which in turn requires rejecting the underlying quantum-level agreement structure.
\LH{IA17}
\end{corollary}
\leanmeta{\LH{IA17}}

Therefore: ``all matter in this universe'' is a technical term. It means:
all matter that agrees with us on the invariant set $\{c, \hbar, k_B, \ldots\}$.
This is not a rhetorical claim; it is the \emph{definition} of shared universe membership.
